% ---- References (back matter) ----
% Sources cited in the body as [S1]..[S6] and the Appendix A model cards.
\chapter{References}
\label{chap:references}
\markboth{References}{References}

\noindent The body text cites sources as \textsc{[S1]}–\textsc{[S7]}. Each entry below gives the work, its venue, and a stable link (an arXiv identifier or a first-party page). These are the \emph{primary} sources for the system-level mechanisms the book describes; hardware figures and prices are additionally flagged in-text for verification against current vendor documentation.

\begin{enumerate}\tightlist
\item \textbf{Continuous batching (Orca).} G.-I. Yu, J. S. Jeong, G.-W. Kim, S. Kim, and B.-G. Chun, ``Orca: A Distributed Serving System for Transformer-Based Generative Models,'' in \emph{Proc.\ 16th USENIX Symposium on Operating Systems Design and Implementation (OSDI '22)}, USENIX Association, 2022, pp.\ 521--538. \texttt{https://www.usenix.org/system/files/osdi22-yu.pdf}
\item \textbf{PagedAttention / vLLM.} W. Kwon, Z. Li, S. Zhuang, Y. Sheng, L. Zheng, C. H. Yu, J. E. Gonzalez, H. Zhang, and I. Stoica, ``Efficient Memory Management for Large Language Model Serving with PagedAttention,'' in \emph{Proc.\ 29th ACM Symposium on Operating Systems Principles (SOSP '23)}, 2023. arXiv:2309.06180. \texttt{https://arxiv.org/abs/2309.06180}
\item \textbf{Prefix caching (SGLang / RadixAttention).} L. Zheng, L. Yin, Z. Xie, C. Sun, J. Huang, C. H. Yu, S. Cao, C. Kozyrakis, I. Stoica, J. E. Gonzalez, C. Barrett, and Y. Sheng, ``SGLang: Efficient Execution of Structured Language Model Programs,'' 2023. arXiv:2312.07104. \texttt{https://arxiv.org/abs/2312.07104}
\item \textbf{P/D disaggregation (Mooncake; see also Splitwise).} Y. Qin et al., ``Mooncake: A KVCache-centric Disaggregated Architecture for LLM Serving,'' 2024. arXiv:2407.00079. \texttt{https://arxiv.org/abs/2407.00079}; and P. Patel, E. Choukse, C. Zhang, A. Shah, \'{I}. Goiri, S. Maleki, and R. Bianchini, ``Splitwise: Efficient Generative LLM Inference Using Phase Splitting,'' in \emph{Proc.\ ISCA}, 2024. arXiv:2311.18677.
\item \textbf{Speculative decoding (DFlash).} J. Chen, Y. Liang, and Z. Liu, ``DFlash: Block Diffusion for Flash Speculative Decoding,'' 2026. arXiv:2602.06036. \texttt{https://arxiv.org/abs/2602.06036}
\item \textbf{FP8 KV quantization (vLLM).} vLLM project, ``The State of FP8 KV-Cache and Attention Quantization in vLLM,'' vLLM blog, 2026-04-22. \texttt{https://vllm.ai/blog/2026-04-22-fp8-kvcache}. (Reports FP8-vs-BF16 \emph{inter-token-latency (ITL) slope ratios}, a decode-\emph{latency} metric; Table 3, ``FP8 slope (\% of BF16),'' gives Llama-3.1-8B at 54\%. It also describes FP8 as halving KV \emph{storage} in the usual tensor-level sense.) \emph{This blog does NOT supply a measured per-token KV \emph{memory} footprint for the 80-layer hidden-8192 full-MHA geometry, and its 54\% is a latency slope, not a storage ratio.} The book therefore does NOT claim a vLLM-measured ``1.42~MB/token'' memory figure; it uses the \emph{raw tensor} FP8 value (1 byte/element $\Rightarrow$ 2 x 80 x 8192 = 1,310,720~B $\approx$ 1.31~MB/token, exactly 50\% of BF16). Any block-allocation / scale-offset / fragmentation overhead is \emph{not} vendor-measured here and must be established by a deployment benchmark before it is used as a capacity coefficient.
\item \textbf{KV cache offloading (vLLM).} vLLM project, ``Inside vLLM's New KV Offloading Connector,'' vLLM blog, 2026-01-08. \texttt{https://vllm.ai/blog/2026-01-08-kv-offloading-connector}. (This is a \emph{distinct} blog from the FP8-quantization one. Its reported TTFT reduction of 2--22$\times$ comes from a particular warm CPU-KV-cache-hit / cache-reuse benchmark for that connector and published model/prompt configuration; it is \emph{not} a generic offloading speedup that applies to a cold cache or to every workload, and memory block layout here depends on the model, attention backend, and block size.) This source is cited as \textsc{[S7]} for the TTFT-offloading claim, separately from the FP8-KV blog \textsc{[S6]}.
\end{enumerate}

\medskip
\noindent\textbf{Primary sources for the major mechanisms in Chapters 1--26.} Beyond the seven system-level \textsc{[S$n$]} tags cited in the body, the following carry the mechanisms the handbook describes (transformer architecture, the roofline, attention kernels, parallelism, quantization, mixture-of-experts, and benchmarking). They are the provenance layer for the corresponding mechanism discussion.

\begin{itemize}\tightlist
\item \textbf{The transformer.} A. Vaswani et al., ``Attention Is All You Need,'' in \emph{Advances in Neural Information Processing Systems (NeurIPS) 30}, 2017. arXiv:1706.03762. \texttt{https://arxiv.org/abs/1706.03762}
\item \textbf{The roofline model.} S. Williams, A. Waterman, and D. Patterson, ``Roofline: An Insightful Visual Perspective on Peak, Memory Bandwidth, and Arithmetic Intensity,'' \emph{Communications of the ACM} 52(4):65--76, 2009. \texttt{https://doi.org/10.1145/1498765.1498785}
\item \textbf{FlashAttention.} T. Dao, D. Y. Fu, S. Ermon, A. Rudra, and C. R\'e, ``FlashAttention: Fast and Memory-Efficient Exact Attention with IO-Awareness,'' in \emph{NeurIPS 35}, 2022. arXiv:2205.14135; and T. Dao, ``FlashAttention-2: Faster Attention with Better Parallelism and Work Partitioning,'' 2023. arXiv:2307.08691.
\item \textbf{Grouped-query attention (GQA).} J. Ainslie, J. Lee-Thorp, M. de Jong, Y. Zemlyanskiy, F. Lebron, and S. Sanghai, ``GQA: Training Generalized Multi-Query Transformer Models from Multi-Head Checkpoints,'' in \emph{EMNLP}, 2023. arXiv:2305.13245.
\item \textbf{Mixture-of-experts (sparse).} N. Shazeer et al., ``Outrageously Large Neural Networks: The Sparsely-Gated Mixture-of-Experts Layer,'' in \emph{ICLR}, 2017. arXiv:1701.06538; W. Fedus, B. Zoph, and N. Shazeer, ``Switch Transformers: Scaling to Trillion Parameter Models with Simple and Efficient Sparsity,'' \emph{JMLR} 23, 2022. arXiv:2101.03961; and A. Q. Jiang et al., ``Mixtral of Experts,'' 2024. arXiv:2401.04088.
\item \textbf{Tensor and pipeline parallelism.} M. Shoeybi et al., ``Megatron-LM: Training Multi-Billion Parameter Language Models Using Model Parallelism,'' 2019. arXiv:1909.08053; and Y. Huang et al., ``GPipe: Efficient Training of Giant Neural Networks using Pipeline Parallelism,'' in \emph{NeurIPS 32}, 2019. arXiv:1811.06965.
\item \textbf{Quantization.} B. Jacob et al., ``Quantization and Training of Neural Networks for Efficient Integer-Arithmetic-Only Inference,'' in \emph{Proc.\ IEEE/CVF Conference on Computer Vision and Pattern Recognition (CVPR)}, 2018, pp.~2704--2713, arXiv:1712.05877. (Leading author is Benoit Jacob, not the author of the binary-network paper of the same period; \texttt{https://openaccess.thecvf.com/content\_cvpr\_2018/html/Jacob\_Quantization\_and\_Training\_CVPR\_2018\_paper.html}.) Also E. Frantar et al., ``GPTQ: Accurate Post-Training Quantization for Generative Pre-trained Transformers,'' in \emph{ICLR}, 2023. arXiv:2210.17323; and T. Dettmers, A. Pagnoni, A. Holtzman, and L. Zettlemoyer, ``QLoRA: Efficient Finetuning of Quantized LLMs,'' in \emph{NeurIPS 36}, 2023. arXiv:2305.14314.
\item \textbf{Speculative decoding (canonical).} Y. Leviathan, M. Kalman, and Y. Matias, ``Fast Inference from Transformers via Speculative Decoding,'' in \emph{ICML}, 2023. arXiv:2211.17192; and C. Chen et al., ``Accelerating Large Language Model Decoding with Speculative Sampling,'' 2023. arXiv:2302.01318.
\item \textbf{Benchmarking a model, not a leaderboard.} P. Liang et al., ``Holistic Evaluation of Language Models (HELM),'' \emph{TMLR}, 2023. arXiv:2211.09110; and L. Zheng et al., ``Judging LLM-as-a-Judge with MT-Bench and Chatbot Arena,'' in \emph{NeurIPS 36}, 2023. arXiv:2306.05685.
\end{itemize}

\noindent\textbf{Appendix A first-party model cards and reports.} The frontier-model figures in Appendix A are drawn from the vendors' own technical reports and model cards (all \textsc{[1P]}): DeepSeek-V4 (report and card), Kimi K3 (report and card), Qwen3.8-Flash-Next (report and card), and GLM-5.3-Flash (report and card). The links are listed at the end of Appendix A and must be re-opened against the current versions, as the appendix's content is a dated snapshot.

\noindent\textbf{Other provenance.} Home-lab measurements (where the book marks that it measured, rather than sourced) are reproducible from the build scripts in the accompanying repository. General knowledge of inference, memory, and attention behavior not tied to a specific system claim does not carry an \textsc{[S$n$]} tag.
